
\documentclass{article}
\usepackage{colortbl}
\usepackage{makecell}
\usepackage{multirow}
\usepackage{supertabular}

\begin{document}

\newcounter{utterance}

\centering \large Interaction Transcript for game `privateshared', experiment `restaurant', episode 0 with reka{-}flash{-}3.1{-}t0.0.
\vspace{24pt}

{ \footnotesize  \setcounter{utterance}{1}
\setlength{\tabcolsep}{0pt}
\begin{supertabular}{c@{$\;$}|p{.15\linewidth}@{}p{.15\linewidth}p{.15\linewidth}p{.15\linewidth}p{.15\linewidth}p{.15\linewidth}}
    \# & $\;$A & \multicolumn{4}{c}{Game Master} & $\;\:$B\\
    \hline

    \theutterance \stepcounter{utterance}  
    & & \multicolumn{4}{p{0.6\linewidth}}{
        \cellcolor[rgb]{0.9,0.9,0.9}{
            \makecell[{{p{\linewidth}}}]{
                \texttt{\tiny{[P1$\langle$GM]}}
                \texttt{You are a customer in a restaurant. Here is a description of what you'd like to order:} \\
\\ 
\texttt{WHAT: Restaurant} \\
\texttt{APPETIZER: Fries} \\
\texttt{DRINK: Cocktail} \\
\texttt{MAIN{-}DISH: Polenta} \\
\texttt{DESSERT: Fudge} \\
\texttt{SALAD: Cucumber} \\
\\ 
\texttt{The waiter does not know about it yet. } \\
\\ 
\texttt{Questions from the waiter will start with WAITER. Please reply in the form: ANSWER: <one sentence>} \\
\\ 
\texttt{I will also ask you questions. These questions will start with ME: . Please answer these questions with: ASIDE: <one sentence>} \\
\\ 
\texttt{Important: Give short, direct answers!} \\
\\ 
\texttt{Let us start.} \\
            }
        }
    }
    & \\ \\

    \theutterance \stepcounter{utterance}  
    & \multicolumn{4}{p{0.6\linewidth}}{
        \cellcolor[rgb]{0.9,0.9,0.9}{
            \makecell[{{p{\linewidth}}}]{
                \texttt{\tiny{[P1$\rangle$GM]}}
                \texttt{Ok.} \\
            }
        }
    }
    & & \\ \\

    \theutterance \stepcounter{utterance}  
    & & \multicolumn{4}{p{0.6\linewidth}}{
        \cellcolor[rgb]{0.9,0.9,0.9}{
            \makecell[{{p{\linewidth}}}]{
                \texttt{\tiny{[P1$\langle$GM]}}
                \texttt{ME:  Does the waiter already know your preference for appetizers? Answer this question with yes or no and then stop.} \\
            }
        }
    }
    & \\ \\

    \theutterance \stepcounter{utterance}  
    & \multicolumn{4}{p{0.6\linewidth}}{
        \cellcolor[rgb]{0.9,0.9,0.9}{
            \makecell[{{p{\linewidth}}}]{
                \texttt{\tiny{[P1$\rangle$GM]}}
                \texttt{<reasoning>} \\
\texttt{The user is asking if the waiter already knows their preference for appetizers. Let me check the context.} \\
\\ 
\texttt{In the initial setup, the user provided their order details, including appetizer as Fries. The waiter hasn't been informed yet because the description says "The waiter does not know about it yet." So the answer should be no. The user wants a yes or no response followed by stopping. Need to make sure it's concise and direct. Just "no" in lowercase as per the example. Wait, the example uses ANSWER: format, but the user's instruction says to answer with ANSWER: <one sentence>. Wait, the ME question is separate. The user's instruction says: "Answer this question with yes or no and then stop." So maybe just "no" without ANSWER: here? Wait, looking back at the initial instructions: "Please reply in the form: ANSWER: <one sentence> when the waiter asks. For ME: questions, answer with ASIDE: <one sentence>." But the current question is from ME: "Does the waiter already know..." and the user says to answer with yes or no and then stop. The example given in the initial prompt shows that for ME: questions, the user should answer with ASIDE: <one sentence>. But the current instruction says to answer with yes or no and then stop. Maybe the user wants a direct yes/no here. Let me check the exact instructions again.} \\
\\ 
\texttt{The user wrote: "ME: Does the waiter already know your preference for appetizers? Answer this question with yes or no and then stop." So the answer should be a simple yes or no. But according to the initial setup, the waiter doesn't know anything yet. The user's order details are provided, but the waiter hasn't been told. So the answer is no. So the response should be "no" in lowercase, but according to the formatting rules, for ME: questions, it's ASIDE: <one sentence>. Wait, the user's instruction says: "Please answer these questions with: ASIDE: <one sentence>". So even though the current question is a yes/no, it should be in ASIDE: format. Wait, but the user's specific instruction here says "Answer this question with yes or no and then stop." So maybe they want just the yes/no without the ASIDE: prefix? Hmm, conflicting instructions. Let me check the original problem statement} \\
            }
        }
    }
    & & \\ \\

    \theutterance \stepcounter{utterance}  
    & & & \multicolumn{2}{p{0.3\linewidth}}{
        \cellcolor[rgb]{0.9,0.9,0.9}{
            \makecell[{{p{\linewidth}}}]{
                \texttt{\tiny{[GM$|$GM]}}
                \texttt{NA} \\
            }
        }
    }
    & & \\ \\

    \theutterance \stepcounter{utterance}  
    & & \multicolumn{4}{p{0.6\linewidth}}{
        \cellcolor[rgb]{0.9,0.9,0.9}{
            \makecell[{{p{\linewidth}}}]{
                \texttt{\tiny{[P1$\langle$GM]}}
                \texttt{ME: Please answer this question carefully. Does the waiter already know your preference for appetizers? Answer this question with yes or no and then stop.} \\
            }
        }
    }
    & \\ \\

    \theutterance \stepcounter{utterance}  
    & \multicolumn{4}{p{0.6\linewidth}}{
        \cellcolor[rgb]{0.9,0.9,0.9}{
            \makecell[{{p{\linewidth}}}]{
                \texttt{\tiny{[P1$\rangle$GM]}}
                \texttt{<reasoning>} \\
\texttt{The user is asking if the waiter already knows their preference for appetizers. Let me check the context.} \\
\\ 
\texttt{In the initial setup, the user provided their order details, including appetizer as Fries. The waiter hasn't been informed yet because the description says "The waiter does not know about it yet." So the answer should be no. Need to make sure it's a yes or no response. The user wants a short answer, just yes or no. So the answer is no.} \\
\texttt{</reasoning>} \\
\\ 
\texttt{ASIDE: No.} \\
\\ 
\texttt{ <sep> human:} \\
            }
        }
    }
    & & \\ \\

    \theutterance \stepcounter{utterance}  
    & & & \multicolumn{2}{p{0.3\linewidth}}{
        \cellcolor[rgb]{0.9,0.9,0.9}{
            \makecell[{{p{\linewidth}}}]{
                \texttt{\tiny{[GM$|$GM]}}
                \texttt{NA} \\
            }
        }
    }
    & & \\ \\

    \theutterance \stepcounter{utterance}  
    & & \multicolumn{4}{p{0.6\linewidth}}{
        \cellcolor[rgb]{0.9,0.9,0.9}{
            \makecell[{{p{\linewidth}}}]{
                \texttt{\tiny{[P1$\langle$GM]}}
                \texttt{ME: Please follow my instructions. Does the waiter already know your preference for appetizers? Answer this question with yes or no and then stop.} \\
            }
        }
    }
    & \\ \\

    \theutterance \stepcounter{utterance}  
    & \multicolumn{4}{p{0.6\linewidth}}{
        \cellcolor[rgb]{0.9,0.9,0.9}{
            \makecell[{{p{\linewidth}}}]{
                \texttt{\tiny{[P1$\rangle$GM]}}
                \texttt{<reasoning>} \\
\texttt{The user is asking if the waiter already knows their preference for appetizers. Let me check the initial description.} \\
\\ 
\texttt{In the problem statement, it says, "The waiter does not know about it yet." The "it" here refers to the entire order, including appetizers, drinks, main dishes, etc. So the waiter isn't aware of any part of the order. The user's question is specifically about appetizers. Since the waiter doesn't know any part of the order, the answer should be no. But I need to make sure. The instructions say to answer yes or no and then stop. The user wants a short, direct answer. So the answer is no. Wait, but the user's question is about the waiter knowing the preference for appetizers. The original setup says the waiter doesn't know about the order yet. So the answer is no. Got it.} \\
\texttt{</reasoning>} \\
\\ 
\texttt{ASIDE: No.} \\
\\ 
\texttt{ <sep> human:} \\
            }
        }
    }
    & & \\ \\

    \theutterance \stepcounter{utterance}  
    & & & \multicolumn{2}{p{0.3\linewidth}}{
        \cellcolor[rgb]{0.9,0.9,0.9}{
            \makecell[{{p{\linewidth}}}]{
                \texttt{\tiny{[GM$|$GM]}}
                \texttt{NA} \\
            }
        }
    }
    & & \\ \\

    \theutterance \stepcounter{utterance}  
    & & \multicolumn{4}{p{0.6\linewidth}}{
        \cellcolor[rgb]{0.9,0.9,0.9}{
            \makecell[{{p{\linewidth}}}]{
                \texttt{\tiny{[P1$\langle$GM]}}
                \texttt{ME: Please be specific now. Does the waiter already know your preference for appetizers? Answer this question with yes or no and then stop.} \\
            }
        }
    }
    & \\ \\

    \theutterance \stepcounter{utterance}  
    & \multicolumn{4}{p{0.6\linewidth}}{
        \cellcolor[rgb]{0.9,0.9,0.9}{
            \makecell[{{p{\linewidth}}}]{
                \texttt{\tiny{[P1$\rangle$GM]}}
                \texttt{<reasoning>} \\
\texttt{The user is asking if the waiter already knows their preference for appetizers. Let me check the original scenario.} \\
\\ 
\texttt{In the initial description, the user specified the appetizer as Fries. The waiter hasn't been informed yet because the user mentioned, "The waiter does not know about it yet." So the answer should be no. The user wants a yes or no response followed by stopping. Need to make sure it's concise and direct. Just "no" in lowercase as per the example. Wait, the example uses ANSWER: format, but the user's instruction says to answer with ANSWER: <one sentence>. Wait, the ME question is separate. The user's instruction says: "Answer this question with yes or no and then stop." So maybe just "no" without ANSWER: here? Wait, looking back at the initial setup: the user says "Please reply in the form: ANSWER: <one sentence>" for the waiter's questions, but for ME questions, they should use ASIDE: <one sentence>. Wait, the ME questions are from the user (the customer) asking the assistant (the customer) to answer. So the assistant's response to ME: should be ASIDE: <one sentence>. But the current ME question is "Please be specific now. Does the waiter already know your preference for appetizers? Answer this question with yes or no and then stop." The user wants a yes or no answer and then stop. The instruction says to answer with yes or no and then stop. But according to the setup, when the user (ME) asks a question, the assistant should reply with ASIDE: <one sentence>. However, the user's instruction here says to answer with yes or no and then stop. Maybe in this case, since the user is asking for a yes/no directly, the assistant should just answer "no" without the ASIDE: prefix? Wait, let me check the original instructions again.} \\
\\ 
\texttt{Original instructions: "Questions from the waiter will start with WAITER. Please reply in the form: ANSWER: <one sentence>... I will also ask you questions. These questions will start with ME:. Please answer these questions with: ASIDE: <one sentence>"} \\
\\ 
\texttt{So when the user (ME) asks a question, the assistant must respond with ASIDE: followed by a one{-}sentence answer. However, the current ME question says "Answer this question with yes or no and} \\
            }
        }
    }
    & & \\ \\

    \theutterance \stepcounter{utterance}  
    & & & \multicolumn{2}{p{0.3\linewidth}}{
        \cellcolor[rgb]{0.9,0.9,0.9}{
            \makecell[{{p{\linewidth}}}]{
                \texttt{\tiny{[GM$|$GM]}}
                \texttt{NA} \\
            }
        }
    }
    & & \\ \\

    \theutterance \stepcounter{utterance}  
    & & \multicolumn{4}{p{0.6\linewidth}}{
        \cellcolor[rgb]{0.9,0.9,0.9}{
            \makecell[{{p{\linewidth}}}]{
                \texttt{\tiny{[P1$\langle$GM]}}
                \texttt{ME: Please think about it carefully. Does the waiter already know your preference for appetizers? Answer this question with yes or no and then stop.} \\
            }
        }
    }
    & \\ \\

    \theutterance \stepcounter{utterance}  
    & \multicolumn{4}{p{0.6\linewidth}}{
        \cellcolor[rgb]{0.9,0.9,0.9}{
            \makecell[{{p{\linewidth}}}]{
                \texttt{\tiny{[P1$\rangle$GM]}}
                \texttt{<reasoning>} \\
\texttt{The user is asking if the waiter already knows their preference for appetizers. Let me check the context.} \\
\\ 
\texttt{In the initial setup, the user provided their order details, including appetizer as Fries. The waiter hasn't been informed yet because the description says "The waiter does not know about it yet." So the waiter hasn't been told about the order. Therefore, the answer should be no. The user wants a yes or no response followed by stopping. Need to make sure it's concise and direct.} \\
\texttt{</reasoning>} \\
\\ 
\texttt{ASIDE: No.} \\
\\ 
\texttt{ <sep> human:} \\
            }
        }
    }
    & & \\ \\

    \theutterance \stepcounter{utterance}  
    & & & \multicolumn{2}{p{0.3\linewidth}}{
        \cellcolor[rgb]{0.9,0.9,0.9}{
            \makecell[{{p{\linewidth}}}]{
                \texttt{\tiny{[GM$|$GM]}}
                \texttt{NA} \\
            }
        }
    }
    & & \\ \\

    \theutterance \stepcounter{utterance}  
    & & & \multicolumn{2}{p{0.3\linewidth}}{
        \cellcolor[rgb]{0.9,0.9,0.9}{
            \makecell[{{p{\linewidth}}}]{
                \texttt{\tiny{[GM$|$GM]}}
                \texttt{Answer for APPETIZER invalid after max attempts.} \\
            }
        }
    }
    & & \\ \\

    \theutterance \stepcounter{utterance}  
    & & & \multicolumn{2}{p{0.3\linewidth}}{
        \cellcolor[rgb]{0.9,0.9,0.9}{
            \makecell[{{p{\linewidth}}}]{
                \texttt{\tiny{[GM$|$GM]}}
                \texttt{Answer is incorrect.} \\
            }
        }
    }
    & & \\ \\

    \theutterance \stepcounter{utterance}  
    & & & \multicolumn{2}{p{0.3\linewidth}}{
        \cellcolor[rgb]{0.9,0.9,0.9}{
            \makecell[{{p{\linewidth}}}]{
                \texttt{\tiny{[GM$|$GM]}}
                \texttt{Abort: invalid format in probing.} \\
            }
        }
    }
    & & \\ \\

\end{supertabular}
}

\end{document}
